%%%%%%%%%%%%%%%%%%%%%%%%%%%%%%%%%%%%%%%%%%%%%%%%%%%%%%%%
%%%%%%                                            %%%%%%
%%%                                                  %%%
%      Modèle de Rapport.                              %
%               Par Matthieu Maury                     %
%                                                      %
%%%                                                  %%%
%%%%%%                                            %%%%%%
%%%%%%%%%%%%%%%%%%%%%%%%%%%%%%%%%%%%%%%%%%%%%%%%%%%%%%%%

\documentclass[11pt, a4paper]{report}

%%%%%%%%%%%%%%%%%%%%%%%%%%%%%%%%%%%%%%%%%%%%%%%%%%%%%%%%
%% Package essentiel
\usepackage[greek,english]{babel}
\usepackage[T1]{fontenc}
\usepackage{ucs}
\usepackage[utf8x]{inputenc}
\usepackage[top=2.6cm,bottom=2.6cm,right=2.1cm,left=2.1cm]{geometry}
\usepackage{float}

%%%%%%%%%%%%%%%%%%%%%%%%%%%%%%%%%%%%%%%%%%%%%%%%%%%%%%%%
%% Package optionnel
\usepackage{enumerate}
\usepackage{graphicx}
\usepackage{tabularx}
\usepackage{setspace}
\usepackage[dvips]{pstricks}
\usepackage{pstricks-add}
\usepackage{color}
\usepackage{xcolor}
\usepackage{epsfig}
\usepackage{pst-grad} % For gradients
\usepackage{pst-plot} % For axes
\usepackage{amsmath}
\usepackage{amsfonts}
\usepackage{amssymb}
\usepackage{amsxtra}
\usepackage{mathrsfs}
\usepackage{framed}
%\usepackage[framed, thmmarks, amsmath]{ntheorem}
\usepackage{verbatim}
\usepackage{moreverb}
\usepackage{fancyhdr}
\usepackage{url}
\usepackage{listings}
%\usepackage{hyperlinks}
\usepackage{lettrine}

%%%%%%%%%%%%%%%%%%%%%%%%%%%%%%%%%%%%%%%%%%%%%%%%%%%%%%%%
%%%%%%                                            %%%%%%
%%         Configuration de la mise en page           %%
%%%%%                                              %%%%%
%%%%%%%%%%%%%%%%%%%%%%%%%%%%%%%%%%%%%%%%%%%%%%%%%%%%%%%%

%%%%%%%%%%%%%%%%%%%%%%%%%%%%%%%%%%%%%%%%%%%%%%%%%%%%%%%%
%% Profondeur du sommaire
\setcounter{secnumdepth}{4}
\setcounter{tocdepth}{4}

%%%%%%%%%%%%%%%%%%%%%%%%%%%%%%%%%%%%%%%%%%%%%%%%%%%%%%%%
%% Configuration des chapitres
\makeatletter
\def\@makechapterhead#1{%
  \vspace*{50\p@}%
  {\parindent \z@ \raggedright \normalfont
    \interlinepenalty\@M
    \Huge \bfseries\thechapter.\quad#1\par\nobreak
    \vskip 20\p@
  }}
\makeatother

%%%%%%%%%%%%%%%%%%%%%%%%%%%%%%%%%%%%%%%%%%%%%%%%%%%%%%%%
%%%%%%                                            %%%%%%
%%                Début du Document                   %%
%%%%%                                              %%%%%
%%%%%%%%%%%%%%%%%%%%%%%%%%%%%%%%%%%%%%%%%%%%%%%%%%%%%%%%


\lstset{tabsize=3, inputencoding=utf8x, extendedchars=\true, language=C}

%% Symboles des ensembles
\newcommand{\R}{\ensuremath{\mathbb{R}} }
\newcommand{\N}{\ensuremath{\mathbb{N}} }
\newcommand{\Z}{\ensuremath{\mathbb{Z}} }
\newcommand{\Q}{\ensuremath{\mathbb{Q}} }
\newcommand{\C}{\ensuremath{\mathbb{C}} }
\newcommand{\U}{\ensuremath{\mathbb{U}} }
\newcommand{\K}{\ensuremath{\mathbb{K}} }

%% Symboles mathématique
\newcommand{\spi}{\ensuremath{\Pi} } %Pi
\newcommand{\sqqs}{\ensuremath{\forall} } %quelque soit
\newcommand{\sex}{\ensuremath{\exists} } %il existe
\newcommand{\snex}{\ensuremath{\nexists} }
\newcommand{\simpld}{\ensuremath{\Rightarrow} } %implique vers la droite
\newcommand{\simplg}{\ensuremath{\Leftarrow} } %implque vers la gauche
\newcommand{\sequ}{\ensuremath{\Leftrightarrow} }
\newcommand{\sand}{\ensuremath{\wedge} }
\newcommand{\sou}{\ensuremath{\vee} }
\newcommand{\smneg}{\ensuremath{\neg} }
\newcommand{\snimpld}{\ensuremath{\nRightarrow} } %implique vers la droite
\newcommand{\snimplg}{\ensuremath{\nLeftarrow} } %implque vers la gauche
\newcommand{\snequ}{\ensuremath{\nLeftrightarrow} }
\newcommand{\sinclut}{\ensuremath{\in} }
\newcommand{\sninclut}{\ensuremath{\notin} }
\newcommand{\sposd}{\ensuremath{\owns} }
\newcommand{\ssups}{\ensuremath{\supset} }
\newcommand{\snsups}{\ensuremath{\nsupset} }
\newcommand{\ssubs}{\ensuremath{\subset} }
\newcommand{\ssubeq}{\ensuremath{\subseteq} }
\newcommand{\ssupeq}{\ensuremath{\supseteq} }
\newcommand{\snsubeq}{\ensuremath{\nsubseteq} }
\newcommand{\snsupeq}{\ensuremath{\nsupseteq} }
\newcommand{\sneg}{\ensuremath{\neq} }
\newcommand{\saprox}{\ensuremath{\approx} }
\newcommand{\ssim}{\ensuremath{\sim} }
\newcommand{\scompl}[2]{\ensuremath{\complement_{#1}{#2}} }
\newcommand{\sunion}{\ensuremath{\cup} }
\newcommand{\sinters}{\ensuremath{\cap} }
\newcommand{\sx}{\ensuremath{\times} }
\newcommand{\svide}{\ensuremath{\emptyset} }
\newcommand{\Pa}{\ensuremath{\mathcal{P}} }
\newcommand{\lra}{\ensuremath{\longrightarrow}}
\newcommand{\rond}{\ensuremath{\circ} }
\newcommand{\frestric}[2]{\ensuremath{#1|_{#2}} }
\newcommand{\sbar}[1]{\ensuremath{\overline{#1}}}
\newcommand{\si}{\ensuremath{\imath}}
\newcommand{\zl}[1]{\ensuremath{\mathscr{#1}}}
\newcommand{\stimes}{\ensuremath{\times}}
\newcommand{\classe}[1]{\ensuremath{\overset{\bullet}{#1}}}


\begin{document}

%%%%%%%%%%%%%%%%%%%%%%%%%%%%%%%%%%%%%%%%%%%%%%%%%%%%%%%%
%% Inclusion de la page de titre
\pagestyle{fancy}
\renewcommand{\sectionmark}[1]{\markright{\thesection\ #1}}
\renewcommand{\footrulewidth}{0pt}
\renewcommand{\headrulewidth}{0pt}
\fancyhead{} % clear all header fields
\fancyfoot{} % clear all footer fields
\fancyfoot[LO,RE]{\textit{University year 2009-2010}}
\fancyfoot[LE,RO]{\textit{Written with \LaTeX}}

\begin{tabularx}{17cm}{Xr}
  \begin{tabular}{ll}
    Yohann Teston & 881003-P792\\
    \url{yohann.teston@free.fr} &\\
	Matthieu Maury & 860928-P210\\
	\url{mayeu.tik@gmail.com} & \\
  \end{tabular} 

  &
  
  \begin{tabular}{r}
    \includegraphics[width=5cm]{pic/logoupp.eps} \\
    \textit{Department of Information Technology} \\
  \end{tabular}
\end{tabularx}

\vspace{6cm}

\begin{center}
  \textbf{ {\Huge Programming of parallel computers}}\\[0.5em]{\huge Assignment 3 - OpenMP}
\end{center}

\begin{center}
  \today
\end{center}


\newpage


\thispagestyle{empty}
%\input{resume}

\renewcommand{\footrulewidth}{0.5pt}
\renewcommand{\headrulewidth}{0.5pt}
\fancyhead{} % clear all header fields
\fancyhead[RE,LO]{Programming of parallel computers - Assignment 3}


\fancyhead[RO,LE]{\rightmark}

\fancyfoot{} % clear all footer fields
\fancyfoot[LO,RE]{Matthieu Maury \& Yohann Teston}
\fancyfoot[LE,RO]{\thepage}

%Redéfinition du style fancy - plain, utilisé pour les pages de nouveau chapitre
%Le style par défaut est un style plain
\fancypagestyle{plain}{
    \fancyhf{}
    \renewcommand{\headrulewidth}{0pt}

    %Définition des headers identiques à une page normale
	 \fancyfoot[LO,RE]{Matthieu Maury \& Yohann Teston}
    \fancyfoot[LE,RO]{\thepage}
}

\tableofcontents

%\listoffigures

%\newpage

%\doublespacing
%\onehalfspacing
\chapter{Problem description}


\chapter{Solution method}

The program takes two arguments, the length of a vector et their number.

\begin{verbatim}
test the arguments

for i in column
    tmp = norm (v, i, row)

	 for k in row
        q[k][i] = v[k][i]/tmp

    for j from i+1 in col
        sigma = scalar_product(q,i,v,j,row)
        for k in row
            v[k][j] -= sigma*q[k][i]
\end{verbatim}

For the direct parallelization (using the \textit{for} directive), we parallelize the second for loop:
\begin{verbatim}
for j from i+1 in col
\end{verbatim}

We chose to parallelize that loop because it is perfectly parallel. Indeed, 
if we consider $i$ as being a constant, we can see that each iteration in the $j$-loop can be made alone. This is not the case with the outer-most loop.Indeed, the $j$-loop uses data that were computed in a previous iteration of the $i$-loop. Therefore, it is not possible to parallelize that loop.\\ 

In the task parallelization, we use \verb+#pragma omp single+ before the first loop:
\begin{verbatim}
for i in column
\end{verbatim}

and we add the task directive \verb+ #pragma omp task private(j, sigma, k, temp_norm)+ inside the body of the $i$-loop.

\chapter{Results and discussion}


\section{Discussion}

\subsection{Parallelization using the \textit{for} directive}

The Openmp data obtained lead to the same conclusion as the last report: the overhead have to be carefully handled to assure a good use of the processor time.\\

Since we had already data for the pthread implementation of gram, we have compared the data of openmp with the one we get during the pthread lab. Taking pthread has a reference we draw the graph \ref{fig:gram_pthread} page \pageref{fig:gram_pthread}. Wich represent the speedup of OpenMP with pthread has reference.

\begin{figure}[ht]
  \begin{center}
         \resizebox{160mm}{!}{\includegraphics{pic/graph_lab2.eps}}
  \end{center}
  \caption{Gram-Schmidt Openmp comparison speedup with pthread}
  \label{fig:gram_pthread}
\end{figure} 

We see that the OpenMP implementation is equivalent to the pthread one.
(if we do not consider matrices of $10\times 10$ wich are not really relevent du to the massive overhead).


\subsection{Parallelization using the {task}

Those results are explained by the overheads related to task creation. Indeed, the master thread create a lot of small tasks for each iteration. Such operation is very time consuming as OpenMP has to create the data environment of the task, put it in the queue and manage the scheduling.  These tasks must also be joined before continuing with the next iteration of the outermost loop. Of course, this synchronization step takes some time, time that is not used efficiently doing some useful computations.\\

This quite big amount of overhead explains why we got those bad results. Indeed, the best speedup achieved is about 2.5, which is not really good.\\

The results' stagnation is due to the fact that passed a certain number of threads, we do not benefit from having more of them because the overhead becomes too important compared to the computation time. So, an increasing speedup cannot be achieved anymore.

\chapter{Conclusion}


\newpage
\setcounter{page}{1}
\pagenumbering{Roman}
\appendix
\chapter{Source code}

\section{Parallelization using the \textit{for} directive}
\label{for}
\lstinputlisting{../gram_schmidt.c}

\section{Parallelization using the \textit{task} directive}
\label{task}
\lstinputlisting{../gram_schmidt_task.c}

\chapter{Results}

The first column gives the size of the array, the second the number of steps, the third the result with one thread and the last one shows the actual result.

\section{Parallelization using the \textit{for} directive}

\verbatiminput{pic/data_gram.dat}

\section{Parallelization using the \textit{task} directive}

\verbatiminput{pic/data_task.dat}


\end{document}
